Les équations de Navier--Stokes décrivent l'écoulement d'un fluide visqueux.
Les bilans de masse et de quantité de mouvement sous
\conceptlink{annexes:forme_conservative}{forme conservative} sont :
\begin{equation}
    \label{eq:navier_stokes}
    \boxed{
        \begin{aligned}
            \frac{\partial \rho}{\partial t}
            + \nabla \cdot (\rho \mathbf{u})
            &= 0, \\
            \frac{\partial (\rho \mathbf{u})}{\partial t}
            + \nabla \cdot \left(
                \rho \mathbf{u} \otimes \mathbf{u} + p \mathbf{I}
            \right)
            &= \nabla \cdot \boldsymbol{\tau} + \rho \mathbf{g}.
        \end{aligned}}
\end{equation}

La masse volumique est notée $\rho$, la vitesse $\mathbf{u}$,
la pression $p$ et l'accélération de la pesanteur $\mathbf{g}$.
Le tenseur identité est noté $\mathbf{I}$, et $\otimes$ désigne
le produit tensoriel.

Le tenseur $\boldsymbol{\tau}$ représente les contraintes visqueuses.
Pour un \conceptlink{annexes:type_fluids}{fluide newtonien isotrope}, il s'écrit :
\begin{equation}
    \label{eq:viscous_stress}
    \boldsymbol{\tau}
    = \mu \left[\nabla \mathbf{u}
        + (\nabla \mathbf{u})^{\mathsf{T}}\right]
    + \lambda (\nabla \cdot \mathbf{u})\mathbf{I},
\end{equation}
où $\mu$ est la viscosité dynamique de cisaillement et $\lambda$
le second coefficient de viscosité.

Les équations d'Euler correspondent au cas où les contraintes visqueuses
sont négligées : $\boldsymbol{\tau} = \mathbf{0}$.

